% New English translation from checked cote 25, site 67 (PDF page 68).
% AI-translated by GPT-6 Astra (OpenAI, ChatGPT Pro mode).
\scanpage{25/67}
\mspage{\unsure{62}}

The incidence prescheme (resp. its inverse image in T) is then
precisely the preschemes of zeros of either homomorphism,
or again of the section $\hand{\phi}^{m}$. We could denote the incidence
prescheme by $H^{(m)}$; for m=1 we recover that of no. 1. If X is
over P, we can put $\underline{Y}^{(m)} = X\times_{P} \underline{H}^{(m)}$, and similarly define
the notation $\underline{Y}^{m}$ when \hand{$\zeta$} is a point of $\unsure{\mathrm{Gr}}^{m}$ with values in an S'.
Thus the \var{\deleted{\unsure{\underline{Y}}}}{}{deletion}{ega25_25_067-v0001} $Y^{m}_{\hand{\zeta}}$ are the “\uline{linear} \uline{sections}” of X \var{\inserted[typed]{over P}}{over P}{insertion}{ega25_25_067-v0002} (or rather of
the $X_{S'}$\var{\inserted[typed]{over $P_{S'}$}}{ over $P_{S'}$}{insertion}{ega25_25_067-v0003}) by the linear subvarieties of codimension m of P (or
rather of the $P_{S'}$).

\var{\margin{$|$}}{\placedmargin{$|$}}{marginal note placed}{ega25_25_067-v0004}I take this opportunity for a\var{\xout{xx}}{}{typewriter strike-out}{ega25_25_067-v0005} self-criticism of notation,
which could appear as early as no. 1. It concerns the arbitrary correspondence,
to indicate an object $\underline{Y}$ corresponding to X, of the letter
Y to X (so that if X becomes Z, we no longer quite know what\var{\xout{x}}{}{typewriter strike-out}{ega25_25_067-v0006}
to choose). This drawback has already led me to a few temporary inconsistencies
of notation. Perhaps the slightly more general context
with an integer m as here suggests a reasonable solution: write
$\underline{X}^{(m)}$ in place of $\underline{Y}^{(m)}$, hence $\underline{X}^{(1)}$ in place of the $\underline{Y}$ of no. 1. In this
way we would indeed have, apart from a few details ... , $\underline{X}^{(m)(m')} = X^{(m+m')}$\hand{.}
I shall try this notation from now on. Of course, the exponent m
itself invites criticism, for it is common practice in algebraic \unsure{geometry}
to use an exponent to denote the dimension of the varieties
involved. But since we have never used this kind
of convention, I think we keep our hands free in this matter.

We see immediately from the preceding construction of $\underline{X}^{(m)}$ that
there is a canonical\var{\inserted[typed]{P-}}{ P-}{insertion}{ega25_25_067-v0007}isomorphism
\[ X^{(m)} = \underset{1}{\underline{\mathrm{Grass}}}_{m}\var{\xout{\$}}{}{typewriter strike-out}{ega25_25_067-v0008}(\overset{\hand{\vee}}{\underline{F}}) , \]
where $F_{X}$\var{\xout{\unsure{is a}}}{}{typewriter strike-out}{ega25_25_067-v0009} ( $\hand{\simeq} f^{*}(\hand{\underline{\Omega}}_{P/S})(1)$) is the kernel of $\underline{E}_{X} \hand{\longrightarrow} O_{X}(1)$\hand{,}
\tnote{25-067-1}{The source's final-letter doubt in géométrie is marked on geometry. Its change from Y notation to X notation is followed at exactly the points where it is proposed; the earlier and later formulas are retained literally. The singular incidence prescheme followed by plural preschemes of zeros also follows the source.}

\mspage{67}
